% GENERATED FILE. Edit canonical lessons, then run npm run generate:books.
% canonical-source-hash: fnv1a64:57e2a578e9dbb2a1
% canonical-lessons: ES-C06-cafe, ES-C06-por-favor, ES-C06-hablar, ES-C06-trabajar, ES-C06-estudiar, ES-C06-espanol, ES-C06-practice

\chapter{The First Verbs}
\label{ch:first-verbs}

This chapter is generated from the canonical micro-lessons used by Language Ladder.
Each section stays independently resumable and preserves the authored prerequisite order.

\section[café]{café — coffee, and a place to drink it}

\label{lesson:ES-C06-cafe}



\begin{quote}
\textbf{Warm-up.} {[}PAUSE 2s{]} Before you can make a polite request, you need one small thing to request. \textbf{Café} gives you a useful noun and lets the written accent do a job you already understand.
\end{quote}

\begin{sounds}
Say \textbf{ca-FÉ}. A vowel-final word would normally stress the second-to-last syllable, so the acute accent records the unexpected final stress.
\end{sounds}

\begin{cousinweb}
Spanish \textbf{café} came through Italian \textbf{caffè} and Turkish \textbf{kahve} from Arabic \textbf{qahwah}. English \textbf{coffee} reached the same travelling word by a different European route. The drink and the word both crossed languages.

\textbf{Café} can name the drink or the place, just as English \textbf{coffee} and \textbf{café} divide those jobs. It is masculine: \textbf{el café}.
\end{cousinweb}

\subsection*{Guided Practice}
{[}PAUSE 1s{]}

\begin{itemize}
\raggedright
  \item {[}YOU SAY: “café” — final stress{]}
  \item {[}YOU SAY: “el café”{]}
  \item {[}YOU SAY: the route — “qahwah $\to$ kahve $\to$ caffè $\to$ café”{]}
\end{itemize}

\begin{tcolorbox}[breakable,colback=teal!4,colframe=teal!35!black,title={Wrap-up recall}]
{[}PAUSE 3s{]} Where does the stress fall? (\textbf{On the final é.}) Why is the accent written? (\textbf{It overrides the usual vowel-final stress.}) What route carried the word into Spanish? (\textbf{Arabic qahwah $\to$ Turkish kahve $\to$ Italian caffè $\to$ Spanish café.})
\end{tcolorbox}
\section[por favor]{por favor — \textquotedblleft{}please,\textquotedblright{} i.e. \textquotedblleft{}for a favour\textquotedblright{}}

\label{lesson:ES-C06-por-favor}



\begin{quote}
\textbf{Warm-up.} {[}PAUSE 2s{]} You can already thank (\emph{gracias}) and reply (\emph{de nada}). The third courtesy completes the set: \textbf{por favor}, \textquotedblleft{}please\textquotedblright{} — literally \emph{\textquotedblleft{}for {[}a{]} favour.\textquotedblright{}} You're asking someone to do you a kindness.
\end{quote}

\begin{sounds}
\begin{itemize}
\raggedright
  \item \texttt{r-tap} — one soft tapped \emph{r} in \emph{por} and \emph{favor}.
  \item \texttt{v-b} — Spanish \emph{v} and \emph{b} sound \textbf{identical} (a soft \emph{b}): \emph{favor} $\approx$ \emph{fah-BOR}. (So \emph{favor} and English \emph{favour} look alike \emph{and} the \emph{v} even sounds like the \emph{b} you'd hear in \textquotedblleft{}harbour.\textquotedblright{})
\end{itemize}
\end{sounds}

\begin{cousinweb}
\begin{itemize}
\raggedright
  \item \textbf{por} = \textquotedblleft{}for / through / by,\textquotedblright{} from Latin \textbf{prō} (\textquotedblleft{}for, on behalf of\textquotedblright{}) — the same \emph{pro-} in English \textbf{pro}noun, \textbf{pro}pose, \textbf{pro}tect.
  \item \textbf{favor} = \textquotedblleft{}a favour, goodwill,\textquotedblright{} from Latin \textbf{favēre}, \emph{\textquotedblleft{}to be favourable to, to support.\textquotedblright{}} It's the very word English borrowed: \textbf{favour}, \textbf{favorite}, \textbf{favorable}, \textbf{favouritism}.
\end{itemize}

So \emph{por favor} literally says \textbf{“for a favour”} — \emph{{[}do this{]} as a kindness to me.} English \textbf{favour}, \textbf{favorite}, and \textbf{favorable} keep the same Latin family visible without asking you to learn another language's formula.
\end{cousinweb}

\begin{grammarlens}[title={Grammar Lens: the courtesy set is complete}]
You now hold the three pillars of politeness:

\begin{itemize}
\raggedright
  \item \textbf{por favor} — please (asking).
  \item \textbf{gracias} — thank you (receiving).
  \item \textbf{de nada} — you're welcome (replying).
\end{itemize}
\end{grammarlens}

\subsection*{Guided Practice}
{[}PAUSE 1s{]}

\begin{itemize}
\raggedright
  \item {[}YOU SAY: \textquotedblleft{}por favor\textquotedblright{} — \emph{por fah-BOR}, soft \emph{v/b}{]}
  \item {[}YOU SAY: \textquotedblleft{}Café, por favor\textquotedblright{} — \textquotedblleft{}Coffee, please\textquotedblright{}{]}
  \item {[}YOU SAY: \textquotedblleft{}favor\textquotedblright{} then English \textquotedblleft{}favour, favorite\textquotedblright{} — one family{]}
\end{itemize}

\begin{tcolorbox}[breakable,colback=teal!4,colframe=teal!35!black,title={Wrap-up recall}]
{[}PAUSE 3s{]} What does \emph{por favor} literally mean? (\textquotedblleft{}For a favour.\textquotedblright{}) What Latin verb is \emph{favor} from, and its English cousins? (\emph{favēre} — favour, favorite, favorable.) Which three-word courtesy loop do you now own? (\emph{Por favor $\to$ gracias $\to$ de nada}.) Next: your first regular \textbf{verb} — \emph{hablar}, “to speak.”
\end{tcolorbox}
\section[hablar]{hablar — \textquotedblleft{}to speak,\textquotedblright{} and the engine of Spanish sentences}

\label{lesson:ES-C06-hablar}



\begin{quote}
\textbf{Warm-up.} {[}PAUSE 2s{]} Fixed phrases got you this far. \textbf{Hablar} (“to speak”) gives you a reusable singular \textbf{-ar} pattern. One ending swap says “I speak,” asks a friend, or addresses someone formally.
\end{quote}

\begin{sounds}
\begin{itemize}
\raggedright
  \item \texttt{silent-h} — the \textbf{h is completely silent}: \emph{hablar} = \emph{ah-BLAR}. (Spanish \emph{h} never makes a sound.)
  \item \texttt{v-b} — the \emph{b} is a soft \emph{b}; \emph{hablo} = \emph{AH-bloh}.
\end{itemize}
\end{sounds}

\begin{cousinweb}
\textbf{hablar} comes from Latin \textbf{fābulārī}, \emph{\textquotedblleft{}to chat, to tell tales\textquotedblright{}} — from \textbf{fābula}, \textquotedblleft{}a story.\textquotedblright{} So to \emph{speak}, in Spanish, is at root to \emph{tell stories}. That \textbf{fābula} is thoroughly English: \textbf{fable}, \textbf{fabulous} (\textquotedblleft{}story-like\textquotedblright{}), \textbf{confabulate}, and \textbf{ineffable} (\textquotedblleft{}un-speak-able\textquotedblright{}).

The initial Latin \textbf{f-} weakened and eventually disappeared in this word, leaving the written but silent Spanish \textbf{h-}. The single path you need today is \textbf{fābulārī $\to$ hablar}; later words will let you test how broadly the same historical pattern reaches.
\end{cousinweb}

\begin{grammarlens}[title={Grammar Lens: the -ar present tense (the big one)}]
\textbf{Hablar} belongs to the \textbf{-ar} family. To conjugate the three singular forms you need now, \textbf{drop -ar} and add the ending for the person:

\noindent
\begin{tabularx}{\linewidth}{@{}>{\raggedright\arraybackslash}X>{\raggedright\arraybackslash}X>{\raggedright\arraybackslash}X>{\raggedright\arraybackslash}X@{}}
\toprule
\textbf{person} & \textbf{ending} & \textbf{\emph{hablar} $\to$} & \textbf{meaning} \\
\midrule
I (the ending carries the subject) & \textbf{-o} & \textbf{hablo} & I speak \\
tú (you) & \textbf{-as} & \textbf{hablas} & you speak \\
usted (you, formal) & \textbf{-a} & \textbf{habla} & you speak \\
\bottomrule
\end{tabularx}

The \textbf{-o} ending already says “I,” so Spanish normally leaves the subject pronoun unstated. That is \textbf{pro-drop}: \emph{hablo} is a complete “I speak.”

\textbf{This same template fits every regular -ar verb.} Next you'll snap two more verbs straight into it.
\end{grammarlens}

\subsection*{Guided Practice}
{[}PAUSE 1s{]}

\begin{itemize}
\raggedright
  \item {[}YOU SAY: \textquotedblleft{}hablar\textquotedblright{} — \emph{ah-BLAR}, silent \emph{h}{]}
  \item {[}YOU SAY: the pattern — “hablo, hablas, habla” (I / tú / usted){]}
  \item {[}YOU SAY: \textquotedblleft{}hablar\textquotedblright{} then English \textquotedblleft{}fable, fabulous, ineffable\textquotedblright{} — the \emph{fābula} root{]}
\end{itemize}

\begin{tcolorbox}[breakable,colback=teal!4,colframe=teal!35!black,title={Wrap-up recall}]
{[}PAUSE 3s{]} What Latin word is \emph{hablar} from, and what English word shares it? (\emph{fābulārī} $\leftarrow$ \emph{fābula} — fable.) What happened to its initial Latin \emph{f-}? (It weakened away, leaving silent Spanish \emph{h-}.) How do you say “I speak,” and why is no subject pronoun required? (\emph{Hablo} — the \emph{-o} ending carries “I.”) Next: \textbf{trabajar}, “to work.”
\end{tcolorbox}
\section[trabajar]{trabajar — \textquotedblleft{}to work,\textquotedblright{} and it once meant \emph{torture}}

\label{lesson:ES-C06-trabajar}



\begin{quote}
\textbf{Warm-up.} {[}PAUSE 2s{]} A second \textbf{-ar} verb — snap it into the exact template you just learned. And \emph{trabajar} has the darkest, best etymology in this chapter: it began as a word for \textbf{torture}, and it's the secret origin of English \textbf{travel}.
\end{quote}

\begin{sounds}
\begin{itemize}
\raggedright
  \item \texttt{j-jota} — the \textbf{j} is a raspy \emph{h}-from-the-throat (the \emph{jota}): \emph{trabajar} = \emph{tra-ba-\textbf{HAR}}.
  \item \texttt{r-tap}, \texttt{b-soft} — tapped \emph{r}, soft \emph{b}.
\end{itemize}
\end{sounds}

\begin{cousinweb}
\textbf{trabajar} comes from Vulgar Latin \textbf{tripaliāre}, \emph{\textquotedblleft{}to torture\textquotedblright{}} — from \textbf{tripalium}, a Roman instrument of torture made of \textbf{three stakes} (\emph{tri-} \textquotedblleft{}three\textquotedblright{} + \emph{pālus} \textquotedblleft{}stake,\textquotedblright{} the root of English \textbf{pale} and \textbf{impale}).

The chain of meaning is grimly logical: \emph{torture} $\to$ \emph{torment / suffering} $\to$ \emph{hard toil} $\to$ just \textbf{work}. And English took the very same word down a parallel road:

\begin{itemize}
\raggedright
  \item \textbf{travail} — painful, laborious effort (and the pains of childbirth).
  \item \textbf{travel} — yes: journeying was once such an \emph{ordeal} that \textquotedblleft{}travail\textquotedblright{} (toil) became \textbf{travel} (to make a journey). Every time you \emph{travel}, you're etymologically being \emph{tortured}.
\end{itemize}

So even plain \textbf{trabajo}, “I work,” hides the old idea of toil as torment — a dark little joke inside an everyday verb.
\end{cousinweb}

\begin{grammarlens}[title={Grammar Lens: same -ar template}]
Drop \textbf{-ar}, add the ending — identical to \emph{hablar}:

\noindent
\begin{tabularx}{\linewidth}{@{}>{\raggedright\arraybackslash}X>{\raggedright\arraybackslash}X>{\raggedright\arraybackslash}X@{}}
\toprule
\textbf{person} & \textbf{form} & \textbf{meaning} \\
\midrule
I (ending carries the subject) & \textbf{trabajo} & I work \\
tú & \textbf{trabajas} & you work \\
usted & \textbf{trabaja} & you (formal) work \\
\bottomrule
\end{tabularx}

Notice you didn't learn a new ending — that's the point. \textbf{One template, every regular -ar verb.}
\end{grammarlens}

\subsection*{Guided Practice}
{[}PAUSE 1s{]}

\begin{itemize}
\raggedright
  \item {[}YOU SAY: \textquotedblleft{}trabajar\textquotedblright{} — \emph{tra-ba-HAR}, raspy \emph{j}{]}
  \item {[}YOU SAY: \textquotedblleft{}trabajo, trabajas, trabaja\textquotedblright{} — same endings as \emph{hablo/hablas/habla}{]}
  \item {[}YOU SAY: \textquotedblleft{}trabajar\textquotedblright{} $\to$ travail $\to$ travel — torture became a journey{]}
\end{itemize}

\begin{tcolorbox}[breakable,colback=teal!4,colframe=teal!35!black,title={Wrap-up recall}]
{[}PAUSE 3s{]} What did \emph{trabajar} originally mean, and from what device? (\textquotedblleft{}To torture,\textquotedblright{} $\leftarrow$ \emph{tripalium}, a three-stake tool.) What two English words come from the same root? (\emph{Travail} and \emph{travel}.) How do you say \textquotedblleft{}I work\textquotedblright{}? (\emph{Trabajo} — same \emph{-o} ending.) Next: \textbf{estudiar}, \textquotedblleft{}to study.\textquotedblright{}
\end{tcolorbox}
\section[estudiar]{estudiar — \textquotedblleft{}to study,\textquotedblright{} and the eagerness inside it}

\label{lesson:ES-C06-estudiar}



\begin{quote}
\textbf{Warm-up.} {[}PAUSE 2s{]} Your third \textbf{-ar} verb. By now the pattern should feel automatic — that's exactly what we want. \emph{estudiar} also hides a nicer feeling than the classroom suggests: at root it means \textbf{eagerness.}
\end{quote}

\begin{sounds}
\begin{itemize}
\raggedright
  \item \texttt{st-no-e} — Spanish props up \emph{st-} with an \emph{e}: \textbf{es-tudiar} (the same reflex behind \emph{estar}, \emph{España}, \emph{escuela}). \emph{es-too-DYAR}.
  \item \texttt{diphthong-io} — the \emph{-iar} glides: \emph{es-too-\textbf{DYAR}}.
\end{itemize}
\end{sounds}

\begin{cousinweb}
\textbf{estudiar} comes from Latin \textbf{studēre}, \emph{\textquotedblleft{}to be eager, to apply oneself zealously.\textquotedblright{}} A \emph{studium} was \textbf{zeal, enthusiasm, devotion} — only later \textquotedblleft{}study.\textquotedblright{} So to \emph{study} is, at heart, to \emph{throw yourself eagerly} at something.

That \textbf{studēre} is everywhere in English:

\begin{itemize}
\raggedright
  \item \textbf{student} — literally \textquotedblleft{}one who is \emph{eager}.\textquotedblright{}
  \item \textbf{studio}, \textbf{studious}, \textbf{study} — the place, the trait, the act.
  \item Italian \textbf{studio} and the musical \emph{studio} keep the \textquotedblleft{}devoted practice\textquotedblright{} sense.
\end{itemize}

So \emph{Estudio español} is \textquotedblleft{}I \emph{devote myself eagerly} to Spanish\textquotedblright{} — a nicer picture than mere schoolwork.
\end{cousinweb}

\begin{grammarlens}[title={Grammar Lens: the template holds}]
Drop \textbf{-ar}, add the ending — for the third time:

\noindent
\begin{tabularx}{\linewidth}{@{}>{\raggedright\arraybackslash}X>{\raggedright\arraybackslash}X>{\raggedright\arraybackslash}X@{}}
\toprule
\textbf{person} & \textbf{form} & \textbf{meaning} \\
\midrule
I (ending carries the subject) & \textbf{estudio} & I study \\
tú & \textbf{estudias} & you study \\
usted & \textbf{estudia} & you (formal) study \\
\bottomrule
\end{tabularx}

Three verbs, one pattern: \emph{hablo, trabajo, estudio}. The singular regular \textbf{-ar} present is now reusable rather than memorized verb by verb.
\end{grammarlens}

\subsection*{Guided Practice}
{[}PAUSE 1s{]}

\begin{itemize}
\raggedright
  \item {[}YOU SAY: \textquotedblleft{}estudiar\textquotedblright{} — \emph{es-too-DYAR}{]}
  \item {[}YOU SAY: \textquotedblleft{}estudio, estudias, estudia\textquotedblright{} — the now-familiar endings{]}
  \item {[}YOU SAY: \textquotedblleft{}estudiar\textquotedblright{} then English \textquotedblleft{}student, studio\textquotedblright{} — the \emph{studēre} \textquotedblleft{}eagerness\textquotedblright{}{]}
\end{itemize}

\begin{tcolorbox}[breakable,colback=teal!4,colframe=teal!35!black,title={Wrap-up recall}]
{[}PAUSE 3s{]} What did \emph{studēre} originally mean? (\textquotedblleft{}To be eager / apply oneself.\textquotedblright{}) Its English cousins? (Student, studio, studious.) Say \textquotedblleft{}I study\textquotedblright{} — and notice it uses the same \emph{-o} as \emph{hablo} and \emph{trabajo}. (\emph{Estudio}.) Next: put a verb and a noun together into your \textbf{first real sentence}.
\end{tcolorbox}
\section[Hablo español]{Hablo español — your first real sentence}

\label{lesson:ES-C06-espanol}



\begin{quote}
\textbf{Warm-up.} {[}PAUSE 2s{]} This is the moment the course has been building to: you take a \textbf{verb} you conjugated and a \textbf{noun}, and make a sentence \textbf{nobody handed you pre-assembled}. \emph{Hablo español} — \textquotedblleft{}I speak Spanish.\textquotedblright{}
\end{quote}

\begin{cousinweb}
\textbf{español} = \textquotedblleft{}Spanish\textquotedblright{} (the language, and a Spaniard), from \textbf{Hispania}, the Roman name for the Iberian peninsula. The path \emph{Hispania $\to$ hispaniolus $\to$ español} wore the word down, and — look — the \textbf{ñ} is back, exactly the doubled-\emph{n}-under- a-tilde you met in the writing lesson (\emph{Hispaniolus} $\to$ \emph{españ-ol}). Say it \emph{es-pa-\textbf{NYOL}}.

(English \textquotedblleft{}Spanish,\textquotedblright{} \textquotedblleft{}Hispanic,\textquotedblright{} and \textquotedblleft{}Hispania\textquotedblright{} are all the same root — and, fittingly, \emph{español} itself may trace back through Latin to a \textbf{Phoenician} name for the coast, \textquotedblleft{}land of hyraxes / rabbits.\textquotedblright{})
\end{cousinweb}

\begin{grammarlens}[title={What you've built — the sentence assembled}]
Snap two things you already own together:

\begin{quote}
\textbf{Hablo} (I speak, from \emph{hablar}) + \textbf{español} (Spanish) = \textbf{Hablo español.}
\end{quote}

That's it — a complete, grammatical Spanish sentence. Note what's \emph{not} there:

\begin{itemize}
\raggedright
  \item \textbf{No separate subject word.} The \emph{-o} of \emph{hablo} already says “I,” reusing the pro-drop pattern from the previous lesson.
  \item \textbf{No \textquotedblleft{}the\textquotedblright{} or \textquotedblleft{}a.\textquotedblright{}} Languages are bare: \emph{hablo español}, not \textquotedblleft{}\emph{el} español.\textquotedblright{}
\end{itemize}

Now swap the pieces and watch it generate:

\begin{itemize}
\raggedright
  \item \textbf{Estudio español.} — “I study Spanish.”
  \item \textbf{Trabajo.} — “I work.”
  \item \textbf{¿Hablas español?} — “Do you speak Spanish?” (reuse the \textbf{¿ ?} from the writing lesson).
\end{itemize}
\end{grammarlens}

\subsection*{Guided Practice}
{[}PAUSE 1s{]}

\begin{itemize}
\raggedright
  \item {[}YOU SAY: \textquotedblleft{}Hablo español\textquotedblright{} — \emph{AH-bloh es-pa-NYOL}{]}
  \item {[}YOU SAY: swap the verb — “Estudio español”, “Trabajo”{]}
  \item {[}YOU SAY: ask it — \textquotedblleft{}¿Hablas español?\textquotedblright{} (raise the pitch, ¿ opens the question){]}
\end{itemize}

\begin{tcolorbox}[breakable,colback=teal!4,colframe=teal!35!black,title={Wrap-up recall}]
{[}PAUSE 3s{]} Where does \emph{español} come from, and what letter does its ending bring back? (\emph{Hispania}; the \emph{ñ}.) In \emph{Hablo español}, why is there no separate subject word and no “the”? (The \emph{-o} ending carries “I”; languages take no article.) You just built your first sentence from parts — that's the whole game from here. Next: practice.
\end{tcolorbox}
\section[Practice]{Practice — making sentences with -ar verbs}

\label{lesson:ES-C06-practice}



\begin{quote}
\textbf{Warm-up.} {[}PAUSE 2s{]} A milestone chapter: for the first time you're not reciting fixed phrases, you're \textbf{building sentences from a pattern}. Drill the three verbs through the endings until \emph{hablo / hablas / habla} is automatic.
\end{quote}

\begin{grammarlens}[title={Grammar Lens: conjugate on command}]
{[}PAUSE 1s{]} Run each verb through the three singular jobs you know: the ending alone for \textbf{I} $\to$ \textbf{tú} $\to$ \textbf{usted}:

\begin{itemize}
\raggedright
  \item \textbf{hablar}: hablo · hablas · habla
  \item \textbf{trabajar}: trabajo · trabajas · trabaja
  \item \textbf{estudiar}: estudio · estudias · estudia
\end{itemize}

Same three endings (\textbf{-o / -as / -a}) every time — that's the whole regular \textbf{-ar} present, singular.
\end{grammarlens}

\subsection*{The exchange — say something real}
\begin{quote}
— ¿\textbf{Hablas} español? — \textbf{Hablo español.} — ¿\textbf{Estudias} español? — \textbf{Estudio español.} — \textbf{Café, por favor.} — ¡Gracias! — De nada.
\end{quote}

\begin{grammarlens}[title={What you've built this chapter}]
\begin{itemize}
\raggedright
  \item \textbf{por favor} — please ($\leftarrow$ \textquotedblleft{}for a favour\textquotedblright{}), completing gracias / de nada / por favor.
  \item \textbf{the singular regular -ar pattern} — drop \emph{-ar}, add \textbf{-o / -as / -a}: the template behind \emph{hablar, trabajar,} and \emph{estudiar}.
  \item \textbf{hablar} ($\leftarrow$ \emph{fābula}, fable; the f$\to$h law), \textbf{trabajar} ($\leftarrow$ \emph{tripalium}, torture $\to$ travail/travel), \textbf{estudiar} ($\leftarrow$ \emph{studēre}, eagerness $\to$ student).
  \item \textbf{Hablo español} — your first self-assembled sentence (pro-drop, no article).
\end{itemize}
\end{grammarlens}

\subsection*{Guided Practice}
{[}PAUSE 1s{]}

\begin{itemize}
\raggedright
  \item {[}YOU SAY: conjugate all three verbs for I / tú / usted, without stopping{]}
  \item {[}YOU SAY: answer “¿Hablas español?” and “¿Estudias español?” about yourself{]}
  \item {[}YOU SAY: request “Café, por favor” and run the full gracias/de nada loop{]}
\end{itemize}

{[}REPEAT x2{]} Ask and answer “¿Hablas español? — Hablo español.”

\begin{tcolorbox}[breakable,colback=teal!4,colframe=teal!35!black,title={Wrap-up recall}]
{[}PAUSE 3s{]} What are the three singular \emph{-ar} endings? (\emph{-o / -as / -a}.) Give the “I” form of all three verbs. (\emph{Hablo, trabajo, estudio}.) Why does \emph{Hablo español} need no separate subject word? (The \emph{-o} carries “I.”) You can now generate sentences — the course stops handing you phrases and starts handing you \emph{rules}. Next chapter adds the \emph{-er/-ir} verb families.
\end{tcolorbox}
